% Part: first-order-logic
% Chapter: completeness
% Section: compactness

\documentclass[../../../include/open-logic-section]{subfiles}

\begin{document}

\iftag{FOL}
      {\olfileid[hi]{fol}{com}{com}}
      {\olfileid[hi]{pl}{com}{com}}

\olsection{संहतता प्रमेय}

पूर्णता प्रमेय का एक महत्त्वपूर्ण परिणाम संहतता प्रमेय है। संहतता प्रमेय
कहता है कि यदि !!{sentence}s के किसी समुच्चय का प्रत्येक \emph{परिमित}
उपसमुच्चय संतुष्टनीय है, तो पूरा समुच्चय संतुष्टनीय है---भले ही समुच्चय
स्वयं अनंत हो। यह बात तनिक भी स्पष्ट नहीं है। कम-से-कम पहली दृष्टि में ऐसा
कुछ नहीं दिखता जो !!{sentence}s के ऐसे अनंत समुच्चयों की संभावना को रोकता
हो जो विरोधाभासी हों, पर जिनमें विरोधाभास मानो उनकी अनंत संख्या से ही
उत्पन्न होता हो। संहतता प्रमेय कहता है कि ऐसी स्थिति असंभव है: ऐसा कोई
असंतुष्टनीय अनंत !!{sentence}s-समुच्चय नहीं है जिसका प्रत्येक परिमित
उपसमुच्चय संतुष्टनीय हो। पूर्णता प्रमेय की तरह इसका भी निष्पत्ति से जुड़ा
एक रूप है: यदि !!{sentence}s का कोई अनंत समुच्चय किसी बात को निष्पन्न करता
है, तो उसका कोई परिमित उपसमुच्चय पहले ही वह बात निष्पन्न करता है।

\begin{defn}
  !!{formula}s का समुच्चय $\Gamma$ \emph{परिमिततः संतुष्टनीय} तभी और केवल
  तभी है जब प्रत्येक परिमित $\Gamma_0 \subseteq \Gamma$ संतुष्टनीय हो।
\end{defn}

\begin{thm}[संहतता प्रमेय]
\ollabel{thm:compactness}
किन्हीं भी वाक्यों $\Gamma$ और $!A$ के लिए निम्न कथन सत्य हैं:
\begin{enumerate}
  \item $\Gamma \Entails !A$ तभी और केवल तभी जब कोई परिमित $\Gamma_0
    \subseteq \Gamma$ ऐसा हो कि $\Gamma_0 \Entails !A$।
  \item $\Gamma$ संतुष्टनीय तभी और केवल तभी है जब वह परिमिततः संतुष्टनीय
    हो।
\end{enumerate}
\end{thm}

\begin{proof}
हम (2) सिद्ध करते हैं। यदि $\Gamma$ संतुष्टनीय है, तो ऐसी
\iftag{FOL}{!!a{structure}~$\Struct{M}$}{!!a{valuation}~$\pAssign{v}$} है कि
$\iftag{FOL}{\Sat{M}}{\pSat{v}}{!A}$ सभी $!A \in
\Gamma$ के लिए। निस्संदेह, यह
$\iftag{FOL}{\Struct{M}}{\pAssign{v}}$ भी~$\Gamma$ के प्रत्येक परिमित
उपसमुच्चय को संतुष्ट करती है, अतः $\Gamma$ परिमिततः संतुष्टनीय है।

अब मान लें कि $\Gamma$ परिमिततः संतुष्टनीय है। तब प्रत्येक परिमित
उपसमुच्चय~$\Gamma_0 \subseteq \Gamma$ संतुष्टनीय है। सुदृढ़ता से
(\iftag{FOL}{%
  \tagrefs{prfAX/{fol:axd:sou:cor:consistency-soundness},
    prfSC/{fol:seq:sou:cor:consistency-soundness},
    prfND/{fol:ntd:sou:cor:consistency-soundness},
    prfTab/{fol:tab:sou:cor:consistency-soundness}}}{%
  \tagrefs{prfAX/{pl:axd:sou:cor:consistency-soundness},
    prfSC/{pl:seq:sou:cor:consistency-soundness},
    prfND/{pl:ntd:sou:cor:consistency-soundness},
    prfTab/{pl:tab:sou:cor:consistency-soundness}}}),
प्रत्येक परिमित उपसमुच्चय संगत है। तब
\iftag{FOL}{%
  \tagrefs{prfAX/{fol:axd:ptn:prop:proves-compact},
    prfSC/{fol:seq:ptn:prop:proves-compact},
    prfND/{fol:ntd:ptn:prop:proves-compact},
    prfTab/{fol:tab:ptn:prop:proves-compact}}}{%
  \tagrefs{prfAX/{pl:axd:ptn:prop:proves-compact},
    prfSC/{pl:seq:ptn:prop:proves-compact},
    prfND/{pl:ntd:ptn:prop:proves-compact},
    prfTab/{pl:tab:ptn:prop:proves-compact}}} से $\Gamma$ स्वयं संगत होना
चाहिए। पूर्णता (\olref[cth]{thm:completeness}) से, क्योंकि $\Gamma$ संगत
है, वह संतुष्टनीय है।
\end{proof}

\tagprob{FOL}
\begin{prob}
\olref[fol][com][com]{thm:compactness} का (1) सिद्ध कीजिए।
\end{prob}
\tagendprob

\tagprob{notFOL}
\begin{prob}
\olref[pl][com][com]{thm:compactness} का (1) सिद्ध कीजिए।
\end{prob}
\tagendprob

\iftag{FOL}{%
\begin{ex}
प्रत्येक प्रतिरूप~$\Struct{M}$ में, जो किसी सिद्धांत~$\Gamma$ का है, प्रत्येक
पद~$t$ निस्संदेह~$\Domain{M}$ के किसी !!a{element} को निर्दिष्ट करता है।
क्या हम यह गारंटी दे सकते हैं कि~$\Domain{M}$ के प्रत्येक !!{element} को भी
कोई-न-कोई पद निर्दिष्ट करता है? दूसरे शब्दों में, क्या ऐसे
सिद्धांत~$\Gamma$ हैं जिनके सभी प्रतिरूप आच्छादित हों? संहतता प्रमेय दिखाता
है कि यदि $\Gamma$ के अनंत प्रतिरूप हों, तो ऐसा नहीं है। इसे इस प्रकार
देखें: कोई अनंत प्रतिरूप $\Struct{M}$ लें जो~$\Gamma$ का हो, और $c$ को ऐसा
!!a{constant} मानें जो~$\Gamma$ की भाषा में न हो। $\Delta$ को उन सभी वाक्यों
$\eq/[c][t]$ का समुच्चय मानें जिनमें $t$, भाषा~$\Lang{L}$ का कोई पद है,
और यह~$\Gamma$ की भाषा है; अर्थात्
  \[
  \Delta = \Setabs{\eq/[c][t]}{t \in \Trm[L]}.
  \]
$\Gamma \cup \Delta$ के किसी परिमित उपसमुच्चय को $\Gamma'
\cup \Delta'$ के रूप में लिखा जा सकता है, जहाँ $\Gamma' \subseteq \Gamma$
और $\Delta' \subseteq
\Delta$। चूँकि $\Delta'$ परिमित है, उसमें केवल परिमिततः अनेक पद हो सकते
हैं। $a \in \Domain{M}$ को~$\Domain{M}$ का ऐसा !!a{element} मानें जिसे उनमें
से कोई निर्दिष्ट नहीं करता, और $\Struct{M'}$ को ऐसी !!{structure} मानें जो
$\Struct{M}$ जैसी ही है, पर जिसमें $\Assign{c}{M'} = a$ भी है। चूँकि
$a \neq \Value{t}{M}$ उन सभी~$t$ के लिए जो~$\Delta'$ में आते हैं,
$\Sat{M'}{\Delta'}$। चूँकि $\Sat{M}{\Gamma}$, $\Gamma' \subseteq
\Gamma$, और $c$,~$\Gamma$ में नहीं आता, इसलिए
$\Sat{M'}{\Gamma'}$ भी। कुल मिलाकर, $\Sat{M'}{\Gamma' \cup \Delta'}$
प्रत्येक परिमित उपसमुच्चय $\Gamma' \cup \Delta'$ के लिए, जो
$\Gamma \cup \Delta$ का है। अतः $\Gamma \cup \Delta$ का प्रत्येक
परिमित उपसमुच्चय संतुष्टनीय है। संहतता से $\Gamma \cup \Delta$ स्वयं
संतुष्टनीय है। इसलिए ऐसे प्रतिरूप~$\Sat{M}{\Gamma \cup \Delta}$ हैं। ऐसी
प्रत्येक $\Struct{M}$,~$\Gamma$ का प्रतिरूप है, पर आच्छादित नहीं है, क्योंकि
$\Value{c}{M} \neq
\Value{t}{M}$ सभी पदों~$t$ के लिए जो~$\Lang{L}$ के हैं।
\end{ex}

\begin{ex}
ऐसी !!a{language} $\Lang{L}$ पर विचार करें जिसमें !!{predicate}~$<$,
!!{constant}s $\Obj{0}$, $\Obj{1}$, और !!{function}s $+$, $\times$, तथा
$-$ हों। $\Gamma$ को उन सभी !!{sentence}s का समुच्चय मानें जो इस
!!{language} की !!{structure}~$\Struct{Q}$ में सत्य हैं, जिसका प्रांत~$\Rat$
है और जिसकी व्याख्याएँ स्वाभाविक हैं। $\Gamma$,~$\Lang{L}$ के उन सभी !!{sentence}s का समुच्चय है
जो परिमेय संख्याओं के बारे में सत्य हैं। निस्संदेह,~$\Rat$ में (और~$\Real$
में भी) ऐसी कोई संख्या~$r$ नहीं है जो~$0$ से बड़ी, पर $1/k$ से छोटी हो,
सभी $k \in \PosInt$ के लिए। यदि ऐसी संख्या होती, तो वह
\emph{अतिसूक्ष्म} होती: शून्येतर, पर अनंततः छोटी। संहतता प्रमेय से दिखाया
जा सकता है कि~$\Gamma$ के ऐसे प्रतिरूप हैं जिनमें अतिसूक्ष्म संख्याएँ हैं।
हमारी भाषा में भाग का कोई !!a{function} नहीं है (शून्य से भाग अपरिभाषित है,
और !!{function}s की व्याख्या संपूर्ण फलनों से होनी चाहिए)। फिर भी हम
$r < 1/k$ व्यक्त कर सकते हैं, क्योंकि यह तभी और केवल तभी है जब
$r \cdot k < 1$। अब $c$ को नया !!a{constant} मानें और $\Delta$ रखें:
\[
\{\Obj{0} < c\}
\cup \Setabs{ c \times \num k < \Obj{1} }{k \in \PosInt}
\]
(जहाँ $\num{k} = (\Obj{1} + (\Obj{1} + \dots + (\Obj{1} +
\Obj{1})\dots))$, जिसमें $k$ बार~$\Obj{1}$ है)। किसी परिमित
उपसमुच्चय~$\Delta_0$ के लिए, जो~$\Delta$ का है, कोई~$K$ ऐसा है कि सभी
!!{sentence}s $c \times \num{k} < \Obj{1}$ जो~$\Delta_0$ में हैं, उनके
लिए $k < K$। यदि हम
$\Struct{Q}$ को~$\Struct{Q'}$ तक इस प्रकार विस्तारित करें कि
$\Assign{c}{Q'} = 1/K$, तो $\Sat{Q'}{\Gamma_0 \cup \Delta_0}$ किसी भी
परिमित $\Gamma_0 \subseteq \Gamma$ के लिए, और इसलिए $\Gamma \cup \Delta$
परिमिततः संतुष्टनीय है (अभ्यास: इसे विस्तार से सिद्ध कीजिए)। संहतता से
$\Gamma
\cup \Delta$ संतुष्टनीय है। किसी भी प्रतिरूप~$\Struct{S}$ में, जो
$\Gamma \cup
\Delta$ का है, एक अतिसूक्ष्म संख्या है,
अर्थात्~$\Assign{c}{S}$।
\end{ex}
}{}

\tagprob{FOL}
\begin{prob}
अंकगणित के मानक प्रतिरूप~$\Struct{N}$ में ऐसा कोई !!{element}~$k \in \Domain{N}$
नहीं है जो प्रत्येक सूत्र~$\num{n}
< x$ को संतुष्ट करता हो (जहाँ $\num{n}$ वह $\Obj{0}^{\prime\dots\prime}$ है
जिसमें $n$ बार $\prime$ आता है)। संहतता प्रमेय का उपयोग करके दिखाइए कि
अंकगणित की भाषा के वे !!{sentence}s जो मानक प्रतिरूप~$\Struct{N}$ में सत्य
हैं, ऐसी !!a{structure}~$\Struct{N'}$ में भी सत्य हैं जिसमें ऐसा
!!a{element} है जो प्रत्येक सूत्र $\num{n} < x$ को \emph{वास्तव में}
संतुष्ट करता है।
\end{prob}
\tagendprob

\iftag{FOL}{
\begin{ex}
हम जानते हैं कि !!{identity} वाला प्रथम-क्रम तर्क व्यक्त कर सकता है कि
!!{domain} का आकार किसी न्यूनतम आकार का होना चाहिए: वाक्य~$!A_{\ge n}$ (जो
कहता है ``कम-से-कम $n$ भिन्न वस्तुएँ हैं'') केवल उन्हीं संरचनाओं में सत्य
है जहाँ $\Domain{M}$ में कम-से-कम~$n$ वस्तुएँ हों। अतः यदि हम
  \[
  \Delta = \Setabs{!A_{\ge n}}{n \ge 1}
  \]
लें, तो $\Delta$ का प्रत्येक प्रतिरूप अनंत होना चाहिए। इस प्रकार हम किसी
सिद्धांत में~$\Delta$ जोड़कर गारंटी दे सकते हैं कि उसके केवल अनंत प्रतिरूप
हों: $\Gamma \cup \Delta$ के प्रतिरूप ठीक~$\Gamma$ के अनंत प्रतिरूप हैं।

अतः प्रथम-क्रम तर्क अनंतता व्यक्त कर सकता है। किंतु संहतता प्रमेय दिखाता है
कि वह परिमितता व्यक्त नहीं कर सकता। मान लें कि वाक्यों के किसी समुच्चय
$\Lambda$ को संतुष्ट करने वाली !!{structure}s ठीक सभी परिमित संरचनाएँ हों। तब
$\Delta \cup \Lambda$ परिमिततः संतुष्टनीय है। क्यों? मान लें कि
$\Delta' \cup \Lambda' \subseteq \Delta \cup
\Lambda$ परिमित है, जहाँ $\Delta' \subseteq \Delta$ और $\Lambda'
\subseteq \Lambda$। $n$ को वह सबसे बड़ी संख्या मानें जिसके लिए $!A_{\ge n}
\in \Delta'$। चूँकि सभी परिमित संरचनाएँ $\Lambda$ को संतुष्ट करती हैं,
इसलिए कोई प्रतिरूप~$\Struct{M}$ ऐसा है जिसमें परिमिततः अनेक किंतु $\ge n$ !!{element}s
हैं। पर तब $\Sat{M}{\Delta' \cup \Lambda'}$। संहतता से
$\Delta \cup
\Lambda$ का कोई अनंत प्रतिरूप है, जो इस धारणा का खंडन करता है कि $\Lambda$
को संतुष्ट करने वाली !!{structure}s केवल परिमित हों।
\end{ex}
}{}

\end{document}
